\section{Implementation of Two-Finger Gripper Optimization}
\label{sec:two_finger_gripper_optimization}

This chapter presents the optimization framework developed for a two-finger gripper with a
spline-based gripping contour. The framework supports both evolutionary multi-objective optimization
and Bayesian multi-objective optimization. The evolutionary algorithms are implemented using the
\texttt{pymoo} library~\cite{pymoo}, while Bayesian optimization is implemented using the
\texttt{Ax} platform and its underlying \texttt{BoTorch} modeling
framework~\cite{olson2025ax,botorch}.

For both approaches, the evaluation of candidate gripper designs is delegated to an external service
through an API. The evaluation service constructs the gripper geometry, voxelizes the gripper and
workpieces, and calculates the grasp-quality objectives using a modified version of Snuggle-Pack\cite{snuggle_pack}.
This separation provides a modular architecture in which the optimization logic and the
domain-specific grasp evaluation remain decoupled.

The same decision variables, parameter bounds, geometry-generation pipeline, workpieces,
orientations, and grasp-quality objectives are used for both optimization approaches. This enables
their performance to be compared under a common evaluation budget and evaluation procedure.

\subsection{Decision Variables and Bounds for Two-Finger Gripper Problem}
\label{sec:two_finger_decision_variables}

A candidate gripper is represented by the decision vector
\begin{equation}
    \mathbf{x}
    =
    \begin{bmatrix}
        d_{\mathrm{gn}} &
        l_{\mathrm{gn}} &
        h_1 &
        h_2 &
        h_3
    \end{bmatrix}^{\mathrm{T}},
\end{equation}
where the variables are transmitted to the evaluation API in the order
\begin{equation}
    \mathbf{x}
    =
    [
        \texttt{depth\_gn},
        \texttt{length\_gn},
        \texttt{h\_p1\_spline},
        \texttt{h\_p2\_spline},
        \texttt{h\_p3\_spline}
    ]^{\mathrm{T}}.
\end{equation}

The decision variables and their bounds are summarized in Table~\ref{tab:two_finger_variables}.

\begin{table}[htbp]
\centering
\caption{Decision variables of the two-finger gripper optimization problem.}
\label{tab:two_finger_variables}
\begin{tabular}{llll}
\toprule
\textbf{Variable} & \textbf{Configuration parameter} & \textbf{Lower bound} & \textbf{Upper bound}
\\
\midrule
$d_{\mathrm{gn}}$ & \texttt{depth\_gn}       & 1 & 25 \\
$l_{\mathrm{gn}}$ & \texttt{length\_gn}      & 1 & 50 \\
$h_1$             & \texttt{h\_p1\_spline}   & 1 & 15 \\
$h_2$             & \texttt{h\_p2\_spline}   & 1 & 15 \\
$h_3$             & \texttt{h\_p3\_spline}   & 1 & 15 \\
\bottomrule
\end{tabular}
\end{table}


\nomenclature{$\mathbf{x}$}{Decision vector of the two-finger gripper}
\nomenclature{$d_{\mathrm{gn}}$}{Depth parameter in the common gripper interface}
\nomenclature{$l_{\mathrm{gn}}$}{Length of the gripping contour}
\nomenclature{$h_1,h_2,h_3$}{Spline contour heights}
All five decision variables are treated as continuous variables. Candidate values generated by
either optimization framework are restricted to the bounded search space
\begin{equation}
    \mathcal{X}
    =
    [1,25]
    \times
    [1,50]
    \times
    [1,15]^3.
\end{equation}
\nomenclature{$\mathcal{X}$}{Bounded decision space of the two-finger gripper}

For evolutionary optimization, the variables are initially sampled within these bounds and
subsequently modified by the operators of the selected algorithm. For Bayesian optimization, Ax
initially evaluates the center of the search space and generates additional initialization
candidates using Sobol sampling. After sufficient observations have been collected, Ax transitions
to model-based candidate generation.

\subsection{Optimization Framework for Two-Finger Gripper Optimization}
\label{sec:two_finger_framework}

The two-finger gripper problem is solved using either an evolutionary multi-objective optimization
algorithm or Bayesian multi-objective optimization. Both methods use the same bounded decision
space, objective definitions, and external evaluation service.

The common optimization configuration contains:
\begin{itemize}
    \item Number of objectives $n_{\text{obj}}$
    \item Decision-variable definitions and bounds
    \item Optimization direction 
    \item Choice of optimization method
    \item Evaluation endpoint URL
\end{itemize}

For evolutionary optimization, the additional configuration parameters include:
\begin{itemize}
    \item Population size $N$
    \item Maximum number of generations $G_{\max}$
    \item Choice of evolutionary algorithm
\end{itemize}

For Bayesian optimization, the additional configuration parameters include:
\begin{itemize}
    \item Maximum number of candidate-generation batches $B_{\max}$
    \item Maximum number of candidates requested per batch $b$
    \item Bayesian generation strategy
\end{itemize}

In the Bayesian implementation, a batch represents a candidate-generation step rather than an
evolutionary generation. The maximum requested number of Bayesian evaluations is therefore
\begin{equation}
    E_{\max}^{\mathrm{BO}} = B_{\max}b.
\end{equation}
The actual number of evaluations may be smaller if the generation strategy returns fewer than $b$
candidates in a batch. \nomenclature{$B_{\max}$}{Maximum number of Bayesian candidate-generation
batches}
\nomenclature{$b$}{Maximum number of Bayesian candidates requested per batch}
\nomenclature{$E_{\max}^{\mathrm{BO}}$}{Maximum requested number of Bayesian objective evaluations}

\subsubsection{Supported Optimization Methods for Two-Finger Gripper Optimization}

The framework supports the evolutionary and Bayesian multi-objective optimization methods listed in
Table~\ref{tab:two_finger_algorithms}.
\begin{table}[htbp]
\centering
\caption{Supported multi-objective optimization methods for the two-finger gripper.}
\label{tab:two_finger_algorithms}

\small
\begin{tabular}{
    @{}
    p{0.20\linewidth}
    p{0.17\linewidth}
    p{0.53\linewidth}
    @{}
}
\toprule
\textbf{Family} & \textbf{Identifier} & \textbf{Method} \\
\midrule
Evolutionary
& NSGA2
& Non-dominated Sorting Genetic Algorithm II~\cite{nsga2} \\

Evolutionary
& NSGA3
& Non-dominated Sorting Genetic Algorithm III~\cite{nsga3} \\

Evolutionary
& MOEAD
& Multi-Objective Evolutionary Algorithm based on Decomposition~\cite{moead} \\

Evolutionary
& CMOPSO
& Competitive Multi-Objective Particle Swarm Optimization \\

Evolutionary
& SMSEMOA
& S-Metric Selection Evolutionary Multi-Objective Algorithm~\cite{smsemoa} \\

Bayesian
& AX-MBM
& Ax modular BoTorch model for multi-objective Bayesian optimization%
~\cite{olson2025ax,botorch} \\
\bottomrule
\end{tabular}
\end{table}

For evolutionary algorithms requiring reference directions, such as NSGA-III and MOEA/D, uniform
reference directions are generated using the method of Das and Dennis~\cite{das_dennis}:
\begin{equation}
    \Lambda =
    \texttt{get\_reference\_directions}
    (\text{``das-dennis''},\;n_{\text{obj}},\;n_{\text{partitions}}).
\end{equation}
\nomenclature{$\Lambda$}{Set of reference directions for decomposition-based algorithms}
\nomenclature{$n_{\text{obj}}$}{Number of objectives}

The Bayesian implementation uses the Ax generation strategy identified at runtime as
\texttt{Center\allowbreak+\allowbreak Sobol\allowbreak+\allowbreak MBM\allowbreak:\allowbreak fast}.
This strategy consists of the following stages:

\begin{enumerate}
    \item Evaluation of the center of the bounded decision space
    \item Sobol quasi-random sampling for initial space exploration
    \item Model-based candidate generation using the Ax modular BoTorch model (MBM)
\end{enumerate}

The Center and Sobol stages do not use a surrogate model or an acquisition function. Once the
generation strategy transitions to the MBM stage, Ax uses the available observations to fit
probabilistic surrogate models and optimize an acquisition function for multi-objective candidate
selection. For continuous multi-objective problems, this commonly involves Gaussian-process
surrogates and a hypervolume-improvement acquisition function such as \texttt{qLogNEHVI}.

The exact surrogate and acquisition-function classes are selected by Ax and can depend on the
installed Ax and BoTorch versions, the number of observations, and the experiment configuration.
Therefore, the active generation stage, surrogate model, and acquisition function should be recorded
from the runtime configuration for reproducibility.

\subsection{Initialization of Two-Finger Gripper Optimization}
\label{sec:two_finger_initialization}

The initialization process depends on whether evolutionary or Bayesian optimization is selected.

\subsubsection{Evolutionary Initialization}

The selected evolutionary algorithm initializes a population of $N$ candidate grippers:
\begin{equation}
    \mathcal{P}^{(0)}
    =
    \left\{
        \mathbf{x}^{(0)}_1,
        \mathbf{x}^{(0)}_2,
        \ldots,
        \mathbf{x}^{(0)}_N
    \right\},
    \qquad
    \mathbf{x}^{(0)}_i\in\mathcal{X}.
\end{equation}

Each candidate is sampled within the variable bounds. The complete population is evaluated through
the external API before the selected evolutionary algorithm performs its first selection and
variation step.
\nomenclature{$N$}{Evolutionary population size}
\nomenclature{$\mathcal{P}^{(0)}$}{Initial evolutionary population}

\subsubsection{Bayesian Initialization}

Bayesian optimization does not maintain an evolutionary population. Instead, it maintains a dataset
containing the completed gripper evaluations:
\begin{equation}
    \mathcal{D}_t
    =
    \left\{
        \left(
            \mathbf{x}_i,
            \mathbf{k}(\mathbf{x}_i)
        \right)
    \right\}_{i=1}^{t},
\end{equation}
where
\begin{equation}
    \mathbf{k}(\mathbf{x})
    =
    \begin{bmatrix}
        k_1(\mathbf{x}) &
        k_2(\mathbf{x}) &
        \cdots &
        k_{n_{\text{obj}}}(\mathbf{x})
    \end{bmatrix}^{\mathrm{T}}
\end{equation}
contains the grasp-quality KPIs returned by the evaluation service.
\nomenclature{$\mathcal{D}_t$}{Dataset of completed Bayesian trials after $t$ evaluations}
\nomenclature{$\mathbf{k}(\mathbf{x})$}{Natural grasp-quality objective vector}

The Ax generation strategy initially evaluates the center of the bounded search space. Additional
initialization candidates are generated using Sobol quasi-random sampling. Sobol sampling provides
broad and approximately space-filling coverage of the five-dimensional search space before a
surrogate model is introduced.

After the initialization requirements of the generation strategy have been satisfied, Ax transitions
to the model-based MBM stage. From this point onward, new gripper designs are selected using
surrogate predictions and predictive uncertainty.

\subsection{Objective Function Evaluation of Two-Finger Gripper via API}
\label{sec:two_finger_api_evaluation}

For every candidate, the decision variables are sent to the evaluation API as an ordered list. The
API inserts these values into a complete gripper configuration:
\begin{equation}
    \mathbf{x}
    \longrightarrow
    \mathbf{c}_{\mathrm{gripper}}(\mathbf{x})
    \longrightarrow
    \mathcal{G}(\mathbf{x}),
\end{equation}
where $\mathbf{c}_{\mathrm{gripper}}$ denotes the gripper configuration and
$\mathcal{G}(\mathbf{x})$ the resulting geometry.

The configuration also contains fixed parameters, including the contour type, basic jaw dimensions,
jaw angle, stroke, and tool-center-point translation. The spline contour is selected by the static
configuration and is not itself an optimization variable. The completed configuration is used to
create the jaw geometry using CadQuery~\cite{cadquery2026} which converts it into the mesh
representation required for evaluation.

\subsubsection{Grasp Evaluation}
\label{sec:grasp_eval}

The generated gripper is evaluated against each configured workpiece. The gripper and workpiece
geometries are then voxelized to obtain three-dimensional occupancy grids.

The grasp evaluation uses a modified version of Snuggle-Pack, a three-dimensional packing algorithm
that combines FFT-based spatial analysis with proximity heuristics~\cite{snuggle_pack}. In the
original method, correlations between voxel grids are used to identify collision-free object
placements. Here, the same principle is adapted to evaluate relative gripper--workpiece placements.

Let $s_{\mathcal{G}}$ and $s_{\mathcal{W}}$ denote the occupancy grids of the gripper and workpiece.
Their spatial correlation produces a collision map:
\begin{equation}
    \zeta_{\mathcal{G},\mathcal{W}}(\mathbf{q})
    =
    \mathcal{F}^{-1}
    \left(
        \overline{\mathcal{F}(s_{\mathcal{G}})}
        \,\mathcal{F}(s_{\mathcal{W}})
    \right)(\mathbf{q}),
\end{equation}
where $\mathbf{q}$ is a relative displacement and $\mathcal{F}$ denotes the Fourier transform.
Values greater than zero indicate overlapping voxels, whereas zero values identify collision-free
relative placements.

The collision map is combined with Gaussian proximity distributions to determine how closely the
fingers surround the workpiece at feasible placements. These calculations are performed using
FFT-based convolution and GPU-accelerated PyTorch operations. The resulting grasp quality is
represented by the KPI \texttt{max\_prox}. If the required opening cannot be achieved within the
available gripper stroke, the KPI is set to zero.

\nomenclature{$s_{\mathcal{G}}$}{Gripper voxel occupancy grid}
\nomenclature{$s_{\mathcal{W}}$}{Workpiece voxel occupancy grid}
\nomenclature{$\zeta_{\mathcal{G},\mathcal{W}}$}{Gripper--workpiece collision map}
\nomenclature{$\mathbf{q}$}{Relative voxel-grid displacement} \nomenclature{$\mathcal{F}$}{Discrete
Fourier transform}

\subsubsection{Objective Evaluation}
\label{sec:two_finger_objective_evaluation}

The evaluation service returns one grasp-quality KPI for each evaluated workpiece--orientation
combination. Since a larger proximity KPI represents a better grasp result, the natural optimization
problem is
\begin{equation}
    \max_{\mathbf{x}\in\mathcal{X}} k_m(\mathbf{x}),
    \qquad
    m=1,\ldots,n_{\text{obj}}.
\end{equation}

The natural objective vector is therefore
\begin{equation}
    \mathbf{k}(\mathbf{x})
    =
    \begin{bmatrix}
        k_1(\mathbf{x}) &
        k_2(\mathbf{x}) &
        \cdots &
        k_{n_{\text{obj}}}(\mathbf{x})
    \end{bmatrix}^{\mathrm{T}}.
\end{equation}
\nomenclature{$k_m(\mathbf{x})$}{Grasp-quality KPI for objective $m$}
\nomenclature{$\mathbf{k}(\mathbf{x})$}{Natural maximization objective vector}

The objective-direction convention is handled differently by the two optimization frameworks. Since
\texttt{pymoo} assumes minimization, the evolutionary algorithms negates the objective
vector
\begin{equation}
    \mathbf{f}_{\mathrm{EA}}(\mathbf{x})
    =
    -\mathbf{k}(\mathbf{x})
    =
    \begin{bmatrix}
        -k_1(\mathbf{x}) &
        -k_2(\mathbf{x}) &
        \cdots &
        -k_{n_{\text{obj}}}(\mathbf{x})
    \end{bmatrix}^{\mathrm{T}}.
\end{equation}
\nomenclature{$\mathbf{f}_{\mathrm{EA}}(\mathbf{x})$}{Minimization-form objective vector supplied to
the evolutionary algorithms}

In the Bayesian implementation, the external evaluation service values are left as is, since
\texttt{Ax} assumes maximization:
\begin{equation}
    \mathbf{f}_{\mathrm{BO}}(\mathbf{x})
    =
    \mathbf{k}(\mathbf{x}).
\end{equation}
\nomenclature{$\mathbf{f}_{\mathrm{BO}}(\mathbf{x})$}{Objective observations supplied to Ax}

The maximization direction is specified in the Ax optimization configuration. Consequently, the
observed KPI values do not need to be negated before they are attached to the corresponding Ax
trial. This also preserves the physical meaning of the KPI values when the final Pareto frontier is
returned.

If $n_W$ workpieces and $n_R$ orientations are evaluated separately, the number of objectives is
\begin{equation}
    n_{\text{obj}}=n_W n_R.
\end{equation}
For the currently configured single orientation, each workpiece contributes one objective:
\begin{equation}
    n_{\text{obj}}=n_W.
\end{equation}

For performance comparisons, the results from both optimization frameworks are converted to the
common natural maximization convention $\mathbf{k}(\mathbf{x})$.

\subsection{Optimization Loop for Two-Finger Gripper}
\label{sec:two_finger_optimization_loop}

The external evaluation pipeline is shared by the evolutionary and Bayesian optimization approaches.
The methods differ primarily in how they use the completed evaluations to generate subsequent
candidates.

For every candidate, the common evaluation steps are:
\begin{enumerate}
    \item Send the candidate decision vector to the external evaluation API.
    \item Map the decision variables to the complete gripper configuration.
    \item Generate the spline-based gripper geometry using CadQuery.
    \item Voxelize the gripper and configured workpieces.
    \item Calculate FFT-based collision maps using the modified Snuggle-Pack procedure.
    \item Calculate one proximity KPI for each workpiece--orientation combination.
    \item Return the KPI vector to the selected optimization framework.
\end{enumerate}

\subsubsection{Evolutionary Optimization Loop}

In every evolutionary generation, the selected algorithm generates a population of candidate
decision vectors. The candidates are evaluated through the common evaluation pipeline, and the
resulting KPI values are negated before being supplied to \texttt{pymoo}.

The population update can be expressed abstractly as
\begin{equation}
    \mathcal{P}^{(g+1)}
    =
    \mathcal{A}
    \left(
        \mathcal{P}^{(g)},
        \mathbf{F}_{\mathrm{EA}}^{(g)}
    \right),
\end{equation}
where $\mathcal{A}$ represents the selected evolutionary algorithm and
$\mathbf{F}_{\mathrm{EA}}^{(g)}$ contains the minimization-form objective vectors evaluated in
generation $g$. \nomenclature{$\mathcal{A}$}{Selected multi-objective evolutionary algorithm}
\nomenclature{$g$}{Evolutionary generation index}
\nomenclature{$\mathbf{F}_{\mathrm{EA}}^{(g)}$}{Evolutionary objective values in generation $g$}

Depending on the selected algorithm, the next population is produced using operations such as
selection, crossover, mutation, decomposition, particle movement, or survival based on hypervolume
contribution.

\subsubsection{Bayesian Optimization Loop}

During the model-based stage, Ax and BoTorch use the completed trial dataset $\mathcal{D}_t$ to fit
probabilistic surrogate models:
\begin{equation}
    \widehat{\mathbf{k}}_t
    =
    \mathcal{M}(\mathcal{D}_t),
\end{equation}
where $\mathcal{M}$ denotes the surrogate-modeling procedure. \nomenclature{$\mathcal{M}$}{Bayesian
surrogate-modeling procedure} \nomenclature{$\widehat{\mathbf{k}}_t$}{Probabilistic surrogate
approximation after $t$ evaluations}

The surrogate models approximate the relationship between the gripper parameters and the
workpiece-specific KPI values. In addition to predicted objective values, they estimate predictive
uncertainty in regions containing few or no observations.

An acquisition function $\alpha(\mathbf{x}\mid\mathcal{D}_t)$ uses these predictions to balance:
\begin{itemize}
    \item Exploitation of gripper designs expected to improve the current Pareto frontier
    \item Exploration of uncertain regions of the bounded decision space
\end{itemize}

For a requested batch size $b$, the next candidate batch is generated as
\begin{equation}
    \mathcal{X}_{t+1}
    =
    \operatorname*{arg\,max}_{
        \substack{
            \mathcal{B}\subset\mathcal{X}\\
            |\mathcal{B}|\leq b
        }
    }
    \alpha(\mathcal{B}\mid\mathcal{D}_t).
\end{equation}
\nomenclature{$\mathcal{X}_{t+1}$}{Batch of gripper candidates generated by Bayesian optimization}
\nomenclature{$\alpha$}{Bayesian acquisition function}

Each generated candidate is evaluated using the external service. The KPI observation is then
attached to the corresponding Ax trial:
\begin{equation}
    \mathcal{D}_{t+1}
    =
    \mathcal{D}_t
    \cup
    \left\{
        \left(
            \mathbf{x}_{t+1},
            \mathbf{k}(\mathbf{x}_{t+1})
        \right)
    \right\}.
\end{equation}

After new trials have been completed, the surrogate models are updated using the expanded dataset,
and another candidate batch is generated. This process continues until the Bayesian evaluation
budget is exhausted.

Although multiple candidates can be requested in one Bayesian batch, the current implementation
evaluates the candidates one after another through the external API. The implementation therefore
performs \emph{sequential batch evaluation}, rather than parallel batch evaluation.
\subsection{Termination Criteria for Two-Finger Gripper Optimization}
\label{sec:two_finger_termination}

The termination criteria depend on the selected optimization approach.

\begin{table}[htbp]
\centering
\caption{Termination criteria for evolutionary and Bayesian optimization.}
\label{tab:two_finger_termination}
\begin{tabular}{llll}
\toprule
\textbf{Approach} & \textbf{Criterion} & \textbf{Symbol} & \textbf{Description} \\
\midrule
Evolutionary & Design-space tolerance & $x_{\text{tol}}$ & Change in decision variables below
threshold \\

Evolutionary & Objective-space tolerance & $f_{\text{tol}}$ & Relative objective improvement below
threshold \\

Evolutionary & Maximum generations & $G_{\max}$ & Maximum number of generations reached \\

Evolutionary & Maximum evaluations & $E_{\max}$ & Maximum number of evaluations reached \\

Evolutionary & Hypervolume & $HV_{\max}$ & Hypervolume progression stagnates  \\

Bayesian & Maximum batches & $B_{\max}$ & Maximum number of candidate batches reached \\

Bayesian & Maximum evaluations & $E_{\max}^{\mathrm{BO}}$ & Bayesian evaluation budget exhausted \\
\bottomrule
\end{tabular}
\end{table}

For evolutionary optimization, the design-space and objective-space tolerances are evaluated over a
sliding window of $w$ generations. \nomenclature{$x_{\text{tol}}$}{Design-space convergence
tolerance} \nomenclature{$f_{\text{tol}}$}{Objective-space convergence tolerance}
\nomenclature{$G_{\max}$}{Maximum number of evolutionary generations}
\nomenclature{$E_{\max}$}{Maximum number of evolutionary evaluations}
\nomenclature{$w$}{Number of generations in the convergence window}

The current Bayesian implementation uses budget-based termination. It performs at most $B_{\max}$
candidate-generation batches and requests up to $b$ candidates in each batch. Therefore,
\begin{equation}
    E_{\max}^{\mathrm{BO}}=B_{\max}b.
\end{equation}



\subsection{Performance Monitoring for Two-Finger Gripper}
\label{sec:two_finger_monitoring}

The hypervolume indicator~\cite{hypervolume} is used to monitor both evolutionary and Bayesian
optimization. Before calculating hypervolume, all objective vectors are represented using the same
natural maximization convention:
\begin{equation}
    \mathbf{k}(\mathbf{x})
    =
    [k_1(\mathbf{x}),\ldots,k_{n_{\text{obj}}}(\mathbf{x})]^{\mathrm{T}}.
\end{equation}

For evolutionary optimization, a callback records the optimization state after each generation. Let
$\mathcal{K}_{\mathrm{EA}}^{(g)}$ denote the non-dominated KPI vectors available after generation
$g$. The corresponding hypervolume is
\begin{equation}
    \mathrm{HV}_{\mathrm{EA}}^{(g)}
    =
    \mathrm{HV}
    \left(
        \mathcal{K}_{\mathrm{EA}}^{(g)},
        \mathbf{r}
    \right).
\end{equation}

For Bayesian optimization, hypervolume can be calculated after every completed trial or after every
completed candidate batch. Let $\mathcal{K}_{\mathrm{BO}}^{(t)}$ denote the non-dominated KPI
vectors obtained from the first $t$ completed Bayesian trials. The corresponding hypervolume is
\begin{equation}
    \mathrm{HV}_{\mathrm{BO}}^{(t)}
    =
    \mathrm{HV}
    \left(
        \mathcal{K}_{\mathrm{BO}}^{(t)},
        \mathbf{r}
    \right).
\end{equation}
\nomenclature{$\mathcal{K}_{\mathrm{EA}}^{(g)}$}{Non-dominated KPI set after evolutionary generation
$g$} \nomenclature{$\mathcal{K}_{\mathrm{BO}}^{(t)}$}{Non-dominated KPI set after $t$ Bayesian
evaluations} \nomenclature{$\mathrm{HV}_{\mathrm{EA}}^{(g)}$}{Evolutionary hypervolume after
generation $g$} \nomenclature{$\mathrm{HV}_{\mathrm{BO}}^{(t)}$}{Bayesian hypervolume after $t$
completed evaluations} \nomenclature{$\mathbf{r}$}{Fixed hypervolume reference point}

Under the maximization convention, the fixed reference point $\mathbf{r}$ must be worse than the
objective values of interest. The same reference point and objective scaling must be used for all
compared methods. An increasing hypervolume generally indicates improvement in convergence toward
the Pareto front, coverage of the objective space, or both.

Evolutionary generations and Bayesian batches can contain different numbers of candidate
evaluations. Therefore, the primary comparison is performed using hypervolume as a function of the
cumulative number of objective evaluations:
\begin{equation}
    \mathrm{HV}=\mathrm{HV}(E),
\end{equation}
where $E$ is the number of completed external API evaluations.
\nomenclature{$E$}{Cumulative number of completed external evaluations}

The framework additionally records total wall-clock time because every function evaluation includes
API communication, CadQuery geometry generation, voxelization, and GPU-based FFT spatial analysis.
Hypervolume as a function of wall-clock time may therefore also be used to compare practical
optimization efficiency.



\subsection{Optional Hyperparameter Optimization for Two-Finger Gripper Problem}
\label{sec:two_finger_hyperparameter_optimization}

For the evolutionary algorithms, the framework optionally supports automated tuning of algorithm
hyperparameters, such as crossover probability and mutation rate, using a meta-optimization
approach. A \texttt{HyperparameterProblem} is constructed around the selected algorithm and
evaluated over multiple random seeds. The negative mean hypervolume across the runs is used as the
meta-objective:
\begin{equation}
    \min_{\boldsymbol{\theta}}
    \left\{
        -\frac{1}{|\mathcal{Z}|}
        \sum_{s\in\mathcal{Z}}
        \mathrm{HV}(\mathcal{K}_s,\mathbf{r})
    \right\}.
\end{equation}
\nomenclature{$\boldsymbol{\theta}$}{Evolutionary algorithm hyperparameters}
\nomenclature{$\mathcal{Z}$}{Set of random seeds used for meta-optimization}

The best hyperparameter configuration is subsequently applied to the main evolutionary optimization
run.

The current Bayesian implementation uses the generation strategy, surrogate configuration, and
acquisition-function configuration selected by Ax. A separate meta-optimization of Bayesian model
hyperparameters is therefore not currently performed. The user specifies the candidate batch size
and optimization budget, while Ax manages the transition from initialization to model-based
optimization.

For reproducibility, each Bayesian experiment records:
\begin{itemize}
    \item Ax and BoTorch software versions
    \item Random seed, where supported
    \item Generation strategy
    \item Number of initialization trials
    \item Candidate batch size
    \item Total evaluation budget
    \item Active surrogate and acquisition-function classes after entering MBM
\end{itemize}




\subsection{Output of Two-Finger Gripper Optimization}
\label{sec:two_finger_output}

The optimization produces a set of non-dominated gripper configurations. The Pareto-optimal set is
defined as
\begin{equation}
    \mathcal{P}^{*}
    =
    \left\{
        \mathbf{x}\in\mathcal{P}
        \;\middle|\;
        \nexists\,\mathbf{x}'\in\mathcal{P}:
        \mathbf{x}' \prec \mathbf{x}
    \right\},
\end{equation}
where $\mathbf{x}'\prec\mathbf{x}$ means that
\begin{equation}
    f_m(\mathbf{x}')\leq f_m(\mathbf{x})
    \quad
    \forall m
\end{equation}
and
\begin{equation}
    f_j(\mathbf{x}')<f_j(\mathbf{x})
    \quad
    \text{for at least one objective }j.
\end{equation}

The final output includes:
\begin{itemize}
    \item The Pareto-optimal decision vectors
    $\mathbf{x}=[d_{\mathrm{gn}},l_{\mathrm{gn}},h_1,h_2,h_3]^{\mathrm{T}}$
    \item The corresponding objective KPI values
    \item The hypervolume progression
    $\{\mathrm{HV}^{(0)},\mathrm{HV}^{(1)},\ldots,\mathrm{HV}^{(T)}\}$
    \item Total optimization wall-clock time
\end{itemize}

The non-dominated solutions represent the best trade-offs between grasp quality for the configured
workpieces and orientations.

\nomenclature{$\mathcal{P}^{*}$}{Set of Pareto-optimal gripper configurations}